%% =====================================================================================
\section{The distortion schedules}\label{app:schedule}
%% =====================================================================================

Both certified programs evaluate the schedule in the same way, in double precision with round-to-nearest. Given
knots $(q_0,\delta_0),\dots,(q_n,\delta_n)$ and a threshold $P_D$, the value at a prime $p\ge P_D$ is $d_p=\delta_0$
if $\log p\le\log q_0$, $d_p=\delta_n$ if $\log p\ge\log q_n$, and otherwise
\[
  d_p=(1-u)\,\delta_l+u\,\delta_{l+1},\qquad u=\frac{\log p-\log q_l}{\log q_{l+1}-\log q_l},
\]
where $l$ is the first index with $\log p\le\log q_{l+1}$. For $m=\Mzero$, $P_D=17$ and $d_p=0$ for $p<17$; the
$\nKnotsMain$ knots are listed in \cref{tab:knots-main} (and in \path{certificate-14501/schedule/sched_14501.json}). For $m=\MzeroLean$, $P_D=50$, $d_p=0.05$ for $31\le p<50$ and
$d_p=0$ for $p<31$; the knots are listed in \cref{tab:knots}. The tilt is $\hat\nu_p=\mathrm{RU}(1/\mathrm{RD}(1-d_p))$,
and the distortion parameter of the proofs is $\delta_p=1-1/\hat\nu_p$ (\cref{sec:rounding}). In the programs' command
lines the schedules are given by the environment variables \code{KNOTS}, \code{PD}, \code{PSB} and \code{DSB}
(\cref{app:reproduce}).

\begin{table}[ht]
\centering
\caption{The $\nKnotsMain$ knots $(q,\delta)$ of the schedule for $m=\Mzero$.}
\label{tab:knots-main}
\footnotesize
% generated by make_figures.py from certificate-14501/schedule/sched_14501.json -- do not edit
\begin{tabular}{@{}rlrlrlrlrl@{}}
\toprule
$q$ & $\delta$ & $q$ & $\delta$ & $q$ & $\delta$ & $q$ & $\delta$ & $q$ & $\delta$ \\
\midrule
17 & 0.003750 & 19 & 0.005000 & 23 & 0.006250 & 29 & 0.008750 & 31 & 0.036250 \\
37 & 0.036250 & 41 & 0.055000 & 43 & 0.073750 & 47 & 0.072500 & 53 & 0.098150 \\
59 & 0.106800 & 61 & 0.119050 & 67 & 0.123650 & 71 & 0.131380 & 73 & 0.141290 \\
79 & 0.155230 & 83 & 0.160970 & 89 & 0.172510 & 97 & 0.178660 & 101 & 0.182540 \\
103 & 0.188650 & 107 & 0.198220 & 109 & 0.204200 & 113 & 0.206020 & 127 & 0.209837 \\
131 & 0.217630 & 137 & 0.223648 & 139 & 0.230595 & 149 & 0.234931 & 151 & 0.236266 \\
157 & 0.238825 & 163 & 0.241289 & 167 & 0.242881 & 173 & 0.249996 & 179 & 0.259213 \\
181 & 0.261402 & 191 & 0.264497 & 193 & 0.266549 & 197 & 0.269965 & 199 & 0.271361 \\
211 & 0.267464 & 223 & 0.264006 & 227 & 0.266240 & 229 & 0.268576 & 233 & 0.270686 \\
239 & 0.274956 & 241 & 0.279675 & 251 & 0.288000 & 257 & 0.296790 & 263 & 0.298855 \\
269 & 0.299341 & 271 & 0.301179 & 277 & 0.299224 & 281 & 0.300443 & 283 & 0.299810 \\
293 & 0.294209 & 307 & 0.288986 & 311 & 0.288947 & 313 & 0.290166 & 317 & 0.292580 \\
331 & 0.305797 & 337 & 0.309213 & 347 & 0.314773 & 349 & 0.315866 & 353 & 0.318408 \\
359 & 0.319916 & 367 & 0.318618 & 373 & 0.317664 & 379 & 0.316724 & 383 & 0.316106 \\
389 & 0.315191 & 397 & 0.313993 & 401 & 0.313793 & 409 & 0.315719 & 419 & 0.313074 \\
421 & 0.313538 & 431 & 0.310826 & 433 & 0.311277 & 439 & 0.310119 & 443 & 0.311003 \\
449 & 0.317314 & 457 & 0.319036 & 461 & 0.319885 & 463 & 0.322807 & 467 & 0.323646 \\
479 & 0.322280 & 487 & 0.337726 & 491 & 0.352946 & 499 & 0.358381 & 550 & 0.326000 \\
650 & 0.348360 & 800 & 0.342500 & 1{,}000 & 0.365220 & 1{,}200 & 0.363800 & 1{,}600 & 0.373440 \\
2{,}000 & 0.382450 & 2{,}700 & 0.387760 & 3{,}500 & 0.398100 & 5{,}000 & 0.404570 & 8{,}000 & 0.414350 \\
20{,}000 & 0.434850 & 50{,}000 & 0.449200 & 200{,}000 & 0.460000 & & & & \\
\bottomrule
\end{tabular}

\end{table}

\begin{table}[ht]
\centering
\caption{The knots of the schedule for $m=\MzeroLean$ and the corresponding tilts $\nu=1/(1-\delta)$.}
\label{tab:knots}
\small
\begin{tabular}{@{}lrrrrrrrrr@{}}
\toprule
$q$ & $50$ & $80$ & $130$ & $220$ & $400$ & $800$ & $2{,}000$ & $8{,}000$ & $50{,}000$\\
$\delta$ & $0.07$ & $0.1612$ & $0.2191$ & $0.2679$ & $0.3198$ & $0.345$ & $0.3762$ & $0.4031$ & $0.4092$\\
$\nu$ & $1.0753$ & $1.1922$ & $1.2806$ & $1.3659$ & $1.4702$ & $1.5267$ & $1.6031$ & $1.6753$ & $1.6926$\\
\bottomrule
\end{tabular}
\end{table}

%% =====================================================================================
\section{The budgets by prime ranges}\label{app:budget}
%% =====================================================================================

\Cref{tab:budget-main} refines \cref{tab:budget} for the certificate of \cref{thm:main}. For each range it gives the
number of primes, the range of $\delta_p$, the evaluation used (\cref{sec:eval-overview}), the certified total
$\sum c_p$ and the total of the certified first moments $\sum\E[U_p]$; the latter, $\sumMonePXMain$ over $p\le\PX$,
shows how much the distortion gains. For $p>\PX$ the program records the individual costs only at every $500$-th prime;
the row gives the certified total $C$.

\begin{table}[ht]
\centering
\caption{The certified budget for $m=\Mzero$ by prime ranges (from
\code{certificate-14501/output/dump\_14501.txt.gz}). All sums are rounded up in the sixth decimal.}
\label{tab:budget-main}
\small
% generated by make_figures.py from certificate-14501/output -- do not edit
\begin{tabular}{@{}lrllrr@{}}
\toprule
primes & number & $\delta_p$ & evaluation & $\sum c_p$ & $\sum\mathbb{E}[U_p]$ \\
\midrule
$2\le p<17$ & 6 & 0.0000 & first moment & 0.037108 & 0.037108 \\
$17\le p<31$ & 4 & 0.0038--0.0088 & S/B/H & 0.084588 & 0.090176 \\
$31\le p<50$ & 5 & 0.0363--0.0738 & S/B/H & 0.113205 & 0.147744 \\
$50\le p<100$ & 10 & 0.0982--0.1787 & S/B/H & 0.170212 & 0.318257 \\
$100\le p<200$ & 21 & 0.1825--0.2714 & S/B/H & 0.194684 & 0.613946 \\
$200\le p<500$ & 49 & 0.2640--0.3584 & S/B/H & 0.170063 & 1.053353 \\
$500\le p<1{,}000$ & 73 & 0.3277--0.3649 & S/B/H & 0.086675 & 1.038927 \\
$1{,}000\le p<2{,}000$ & 135 & 0.3638--0.3824 & S/B/H & 0.057922 & 1.246231 \\
$2{,}000\le p<7{,}253$ & 624 & 0.3825--0.4123 & S/B/H & 0.055035 & 2.771613 \\
$7{,}253\le p<20{,}000$ & 1{,}335 & 0.4123--0.4348 & $\tau$-split & 0.017663 & 2.570146 \\
$20{,}000\le p<50{,}000$ & 2{,}871 & 0.4349--0.4492 & $\tau$-split & 0.007209 & 2.588585 \\
$50{,}000\le p\le2\cdot 10^5$ & 12{,}851 & 0.4492--0.4600 & $\tau$-split & 0.004111 & 4.267893 \\
\midrule
$p\le 2\cdot 10^5$ & 17{,}984 & & & 0.998469 & 16.743973 \\
$2\cdot 10^5<p\le 10^9$ & 50{,}829{,}550 & 0.4600 & $\tau$-split & 0.001458 & \\
\bottomrule
\end{tabular}

\end{table}

\Cref{tab:budget-fine} does the same for the certificate of \cref{thm:lean}, with ranges between consecutive knots
of its schedule. It also gives the total $\sum b_p$ of the moment bounds of BBMST with the same schedule, and the ratio
$\sum b_p/\sum c_p$. Over $p\le\PX$ the first moments total $\sumMonePX$ and the moment bounds of BBMST $\sumBBPX$.

\begin{table}[ht]
\centering
\caption{The certified budget for $m=\MzeroLean$ between consecutive knots of the schedule (from
\code{logs/dump16k.txt}). All sums are rounded up in the sixth decimal.}
\label{tab:budget-fine}
\small
% generated by make_figures.py from logs/dump16k.txt and logs/cert16k.txt -- do not edit
\begin{tabular}{@{}lrlrrrr@{}}
\toprule
primes & number & $\delta_p$ & $\sum c_p$ & $\sum\mathbb{E}[U_p]$ & $\sum b_p$ & ratio \\
\midrule
$2\le p<31$ & 10 & 0.0000 & 0.120331 & 0.120331 & 0.120331 & 1.00 \\
$31\le p<50$ & 5 & 0.0500 & 0.110242 & 0.141679 & 0.141679 & 1.29 \\
$50\le p<80$ & 7 & 0.0813--0.1588 & 0.126588 & 0.214777 & 0.214777 & 1.70 \\
$80\le p<130$ & 9 & 0.1656--0.2163 & 0.117619 & 0.274202 & 0.270831 & 2.30 \\
$130\le p<220$ & 16 & 0.2198--0.2640 & 0.127195 & 0.441578 & 0.302992 & 2.38 \\
$220\le p<400$ & 31 & 0.2691--0.3191 & 0.123609 & 0.694849 & 0.291946 & 2.36 \\
$400\le p<800$ & 61 & 0.3199--0.3449 & 0.100274 & 0.978090 & 0.253437 & 2.53 \\
$800\le p<2{,}000$ & 164 & 0.3454--0.3762 & 0.082401 & 1.587360 & 0.238184 & 2.89 \\
$2{,}000\le p<8{,}000$ & 704 & 0.3762--0.4031 & 0.056629 & 2.956076 & 0.205359 & 3.63 \\
$8{,}000\le p<50{,}000$ & 4{,}126 & 0.4031--0.4092 & 0.022132 & 4.839575 & 0.113938 & 5.15 \\
$50{,}000\le p\le2\cdot 10^5$ & 12{,}851 & 0.4092 & 0.003860 & 4.128765 & 0.031409 & 8.14 \\
\midrule
$p\le 2\cdot 10^5$ & 17{,}984 & & 0.990874 & 16.377278 & 2.184877 & 2.21 \\
$2\cdot 10^5<p\le 2\cdot 10^8$ & 11{,}060{,}953 & 0.4092 & 0.006511 & & & \\
\bottomrule
\end{tabular}

\end{table}

%% =====================================================================================
\section{Rounding conventions of the checkers}\label{app:rounding}
%% =====================================================================================

This appendix summarizes \path{src/CERT_NOTES.md} and the header of \path{certificate-14501/src/rigx.c}.
Throughout, ``up'' means that the computed double is an upper bound for the exact value of the formula applied to the
computed inputs, and ``down'' a lower bound; the programs are in round-upward mode, and lower bounds use the negation
idiom of \cref{sec:rounding}.

\emph{Build and self-test.} The sources contain \code{\#pragma STDC FENV\_ACCESS ON} and \code{FP\_CONTRACT OFF}
and refuse to compile with \code{-ffast-math}. With GNU \code{gcc}, \code{-frounding-math} is required and the
programs must be built with \code{-O0}: \code{gcc} ignores the \code{FENV\_ACCESS} pragma and at \code{-O1} and above
may reuse a result across \code{fesetround}; Apple clang~21 honours the pragma at \code{-O2}. At start-up
the programs check that directed rounding and the negation idiom work and that no constant was folded in
round-to-nearest.

\emph{First moment.} By \eqref{eq:EU} and \eqref{eq:U-tau}, $\E[U_p(s)]=T_1/(p-1)-U_1$ with
$T_1=\prod_{q<p}(1+\hat\nu_q/(q-1))$ and $U_1=(p-1)^{-1}\sum_{d<N_1}\nu(d)c(d)/d$, the sum running over
$(p-1)$-smooth $d$. The product $T_1$ is rounded up and the subtracted partial sum $U_1$ down.

\emph{Second moment.} $\E[U_p(s)^2]$ is computed by a finite inclusion--exclusion formula of the form
$T/(p-1)^2-2T_1V/(p-1)+S$, where $T=\prod_{q<p}(1+\hat\nu_q(3q-1)/(q-1)^2)$ and $S$, a sum of nonnegative terms
(the programs assert $\hat\nu_q\le q$, which makes them nonnegative), are rounded up, and the subtracted quantities
$T_1$ and $V$ are rounded down (see \code{src/CERT\_NOTES.md} for the formula). The divisor $(p-1)^2$ is rounded down
(\cref{sec:rounding}), and the bound $\E[U_p^2]/(4d(1-d))$ has its denominator rounded down.

\emph{The exact evaluation} (\code{rigx.c}). All dynamic programs (for $T_H$, for $(T,b)$ given a profile, for
$\tau_s$ and for $(e_2,e_3)$) are positive linear maps, evaluated in round-upward mode with upper bounds for the
coefficients $\P(v_q=a)$, so every stored probability is an upper bound; in-place updates process the cells in an
order in which no updated cell is read again. The numbers $c(d)=1-p^{1-j_0(d)}$, $1-\frac1p$ and $t(p-1)$ are rounded
down, so the thresholds $\beta$ of \cref{sec:SBH} and $X$ of \cref{sec:tausplit} are lower bounds, and the hinges
are evaluated at these lower thresholds, which is valid since they are nonincreasing in the threshold. The suffix
sums \eqref{eq:suffix}, the bounds for the pruned subtrees, the lost masses \eqref{eq:lost} and the margin factors are
rounded up.

\emph{The simpler evaluation} (\code{rigcert.c}). The probabilities $\P(s_A)$ are carried as intervals through the
enumeration, and $U_A(s_A)$ is rounded up. The dynamic programs for the laws of $\tau(s_A)$ and $\tau(s_B)$ produce
upper bounds. In $G(k,u)$, the inputs $u$ and $k/(p-1)$ are rounded up and $c=1+(t-u)/a$ is rounded down, which is
safe because $G$ is monotone in each. The suffix sums for $\E[(\tau(s_B)-c)^+]$ are rounded up. The weights
$R_k=\P_{\mathrm{up}}(\tau(s_A)=k)-\P_{\mathrm{lo}}(\text{enumerated},\tau(s_A)=k)$ use the enumerated mass rounded
down. Finally $1-t$ is rounded down. The program enforces $t\le d_p$ exactly.

\emph{Fractional moments.} $q^{-a}$ is formed by repeated division in the required direction, $(a+1)^\theta$ by
repeated multiplication and a directed square root, and the series \eqref{eq:Etheta} is summed with the geometric
bound for its remainder given in \cref{sec:theta}. The constant $c_\theta$ is rounded up: the program uses the
unconstrained maximum $\frac{d}{\theta-1}(u^*)^{-\theta}$ with $u^*$ rounded down, or $1-d$ if the lower bound for $u^*$
exceeds $1$. The products $\E[\tau(s)^\theta]$ are rounded up, and $(p-1)^\theta$ and $1-d$ down.

\emph{Totals and the tail tests.} The total $\eta$ is the sum of the costs, rounded up. For
\cref{prop:bbmst61}, $\mu_{\mathrm{lo}}=\mathrm{RD}(1-\eta_{\mathrm{up}})$ and
$f_{\mathrm{up}}=\mathrm{RU}(T_{\mathrm{up}}/\mu_{\mathrm{lo}})$; the lower bound for $(\log k+\log\log k-3)^2k$ uses lower
bounds for the logarithms from the truncated series $\log y=2\sum_{i\ge0}z^{2i+1}/(2i+1)$, $z=(y-1)/(y+1)$, whose
terms are positive (with $\log2$ from the same series); the programs check $k\ge10$ and that the base is positive.
The logarithm lower bounds were compared with $50$-digit values.

\emph{Audits.} \code{rigcert.c} was derived from a floating-point prototype, which computed in round-to-nearest,
allowed $t$ slightly above $\delta_p$, could round the tilt below $1/(1-\delta_p)$, computed two truncation tails with
cancellation and truncated a series without bounding its remainder; all of this is corrected in the certified program,
which also refuses a debugging option that disables the tilt. A later audit found the rounding gap described in
\cref{sec:rounding}, which has been repaired.

%% =====================================================================================
\section{Reproducing the computations}\label{app:reproduce}
%% =====================================================================================

All paths are relative to the accompanying folder.

\emph{The certificate for \cref{thm:main}.} In \code{certificate-14501/}:
\begin{alltt}\footnotesize
./run_certificate.sh
\end{alltt}
builds \code{src/rigx.c} with \code{cc -O2 -frounding-math -ffp-contract=off} (\code{-O0} under GNU \code{gcc}),
runs \code{rigx \MzeroNum{} 1e9} with
the environment of the certified run (the $\nKnotsMain$ knots in \code{KNOTS}, \code{PD=17 PSB=0 DSB=0}; the accuracy
settings \code{EPS=1e-14 TINY=1e-30 HEFF=1e-16 TS=4194304 THM=1048576 TBUD=4194304}; the margins \code{MARGIN}$=2^{-20}$
and \code{MARGINL}$=$\code{MARGINL2}$=2^{-13}$), and compares the result lines with \code{output/cert\_14501.txt}
(line breaks added):
\begin{alltt}\footnotesize
m=14501 eta_up=0.999926855181 (A=0.234899856304 B=0.763569005991 C=0.001457990070)
  f_up=9.65174e+09 crit_lo=1.57871e+10 SUCCESS(certified)
CERT eta_up=0.99992685518076208 (0x1.fff669aac979p-1) T_up=705974.99523353134
  mu_lo=7.3144819237924708e-05 f_up=9651742974.9486885 crit_lo=15787081835.245779
  k=50847534 log_k_lo=17.744342183707648 loglog_k_lo=2.8760667148961203
\end{alltt}
It takes about $90$ seconds and $1.2$~GB. The cross-check of \cref{tab:xcheck-main} is run by
\path{xcheck/tools/run_xcheck.sh} with the same environment (about $29$ minutes), and \path{validation/run_bf.py} runs
the brute-force comparisons of \cref{sec:tests}.

\emph{The certificate for \cref{thm:lean}.} In \code{src/}:
\begin{alltt}\footnotesize
cc -O2 -frounding-math -ffp-contract=off -o rigcert rigcert.c -lm   # GNU gcc: -O0
K=50:0.07,80:0.1612,130:0.2191,220:0.2679,400:0.3198,800:0.345
K=$K,2000:0.3762,8000:0.4031,50000:0.4092
KNOTS=$K PSB=31 DSB=0.05 NOSTOP=1 \textbackslash
  ./rigcert 16000 50 200000 2e8 47 16 4.5 1000000 20000 0 -1 0.1 0.44 0.5
\end{alltt}
The arguments are, in order: $m$; the prime $50$ from which the knot schedule applies; the largest prime $\PX$
treated by the comparison bound; the largest prime $\Pmax$ treated explicitly; the largest prime $47$ allowed in
$s_A$; the constants $16$ and $4.5$ of the enumeration cut-off $X_p$; the caps $10^6$ and $2\cdot10^4$ of the
divisor-count tables; a verbosity flag; and an unused smooth schedule (ignored when \code{KNOTS} is set).
\code{PSB} and \code{DSB} give $\delta_p=0.05$ for $31\le p<50$; \code{NOSTOP} only disables an early exit. The
expected output (\code{logs/cert16k.txt}; line breaks added) is
\begin{alltt}\footnotesize
m=16000 eta_up=0.997384393974 (A=0.230571210247 B=0.760302397508 C=0.006510785610)
  f_up=1.22745e+08 crit_lo=2.83863e+09 SUCCESS(certified)
CERT eta_up=0.9973843939732947 (0x1.fea92ad34fe9cp-1) T_up=321051.36453405279
  mu_lo=0.0026156060267052972 f_up=122744542.28049764 crit_lo=2838631802.0518227
  k=11078937 log_k_lo=16.220556296052827 loglog_k_lo=2.7862793450197798
CERT stats: max_enumerated=12114518 t_clamped_to_delta=0
\end{alltt}
With \code{DUMP=file} the program writes the certified per-prime costs for $p\le\PX$ (\path{logs/dump16k.txt}).
The run takes about $2.5$ minutes on one core. The commands of the independent recomputation are listed at the end
of \code{logs/xcheck\_final.md} (code in \code{tests/xcheck\_final/}); \code{tests/cert/run\_small.sh} runs the
small-instance checks of \cref{sec:tests}, and the tests of the comparison theorem are the Python scripts in
\code{tests/review/}.

\emph{The results of \cref{sec:further}.} The scripts \path{squarefree/src/reproduce.sh} (about $30$ seconds) and
\path{odd-moduli/reproduce_all.sh} (about $10$ minutes), and the drivers listed in \path{applications/PACKAGE.md},
rerun the certificates of \cref{sec:squarefree,sec:odd,sec:rings}.

\emph{This paper.} In \code{paper/}, \code{make\_figures.py} regenerates \code{fig/} and \code{gen/} from the
certified outputs, and \code{tectonic main.tex} builds the paper.

\needspace{11\baselineskip}
\emph{Lean.} In \code{lean/}:
\begin{alltt}\footnotesize
lake exe cache get                                 # prebuilt Mathlib
lake build MinModulus.Checker2Sound.Final14501     # Theorem \ref{thm:main} (native check: about 10 min)
lake env lean MinModulus/Checker2Sound/Final14501Audit.lean   # axioms of the final theorems
lake build                                         # Theorem \ref{thm:lean} (native check: about 50 s)
lake build MinModulus.Squarefree.Main              # squarefree theorem
lake build MinModulus.Odd.Final                    # Theorem \ref{thm:odd} (native checks: 80 s each)
lake env lean MinModulus/CheckerSound/Audit.lean   # axioms of MinModulus.not_covers
lake env lean MinModulus/Squarefree/Axioms.lean    # axioms of the squarefree theorem
lake env lean MinModulus/Odd/FinalAudit.lean       # axioms of Theorem \ref{thm:odd}
\end{alltt}
